\chapter{F\&O eligibility and the cross-section}\label{ch:elig}
\chapterlead{Short-sale constraints across the whole cash market: what entry into and exit from the F\&O segment do to
price adjustment and drift, whether anomalies are stronger where shorting is hardest, and how drift and price delay vary
across all stocks.}

\section{Entry into and exit from the F\&O segment (H6)}
Entry gives a stock futures, options and formal short-selling eligibility at once; an eligibility exit removes them while
the stock keeps trading in the cash market. This changes the short-sale constraint for many more stocks than the ban
does, and against the whole cash market as a control group.

\paragraph{Design.} Stocks by calendar year, 2015 to 2025. A stock's cohort is the first year fully after its event; the
event year is a transition year. For price delay I use the Callaway and Sant'Anna (2021) estimator with never-treated and
not-yet-treated controls and base-year covariates (log market capitalisation, log traded value, volatility), with a
stock-cluster bootstrap. Drift exists only at announcements, so for drift after bad news I use a stacked
difference-in-differences at the announcement level. The family statistic is a Bonferroni combination of the four
components.

\begin{table}[t]
\caption[Effects of F\&O entry and exit]{Entry speeds price adjustment; drift after bad news does not weaken}\label{tab:h6}
\exsub{Pre-specified components of the H6 test}
\small\begin{tabular}{@{}>{\bfseries\color{accent}}lcccc@{}}
\toprule
\thd{Component} & \thd{Prediction} & \thd{Estimate} & \thd{s.e.} & \thd{One-sided $p$} \\
\midrule
Entry: price delay & $<0$ & $-0.109$ & 0.023 & $1.3\times10^{-6}$ \\
Entry: drift after bad news & $>0$ & $-0.040$ & 0.018 & 0.99 \\
Exit: price delay & $>0$ & $+0.053$ & 0.040 & 0.091 \\
Exit: drift after bad news & $<0$ & $+0.014$ & 0.020 & 0.77 \\
\bottomrule
\end{tabular}
\exsource{Price delay: Callaway and Sant'Anna estimator, stock-cluster bootstrap. Drift: stacked difference-in-differences at the announcement level. Source: author's calculations.}
\end{table}

Prices of stocks that enter the segment respond much faster to market-wide news (median delay 0.24 before entry), and those
of stocks that leave somewhat more slowly. The Sun and Abraham (2021) estimator gives $-0.063$ (s.e.\ 0.019) for entries,
leads three and four years before entry are close to zero, and the result holds without the large November 2024 batch and
without back-cast market capitalisations. The drift components go the other way: after entry the drift that follows bad
news is stronger, not weaker. H6 is supported through price delay and contradicted through drift. Entry requires a stock
to have grown in size and liquidity, which by itself speeds price adjustment; the covariates and the flat leads limit,
but cannot rule out, selection on growth (Figure~\ref{fig:h6es}).

\begin{figure}[t]
\caption[Price delay around F\&O entry and exit]{Price delay falls after entry and rises a little after exit, with no clear trend before either}\label{fig:h6es}
\exsub{Callaway and Sant'Anna estimates of the change in price delay by year relative to the first full year after the event (cells with at least 19 treated stocks)}
\centering\includegraphics[width=\linewidth]{fig_h6_es}
\exsource{Point estimates; the aggregate post-event effects and their standard errors are in Table~\ref{tab:h6}. Source: author's calculations.}
\end{figure}

\section{Anomalies and short-sale constraints (H5)}
Each month I sort stocks within three constraint groups (F\&O stocks; cash-only stocks with SLB trades in the last three
months; cash-only stocks without) into terciles of a composite mispricing score built from momentum, short-term reversal,
maximum daily return, idiosyncratic volatility, illiquidity and turnover, and measure the most overpriced and most
underpriced terciles against the group's value-weighted return. The model predicts that the overpriced side does worse
where shorting is hardest, while the underpriced side does not differ.

Value-weighted, it does not: the short-leg difference between the most constrained group and F\&O stocks is $+0.04\%$ a
month (s.e.\ 0.20). Equal-weighted, anomalies are far stronger among small cash-only stocks, but on both sides (short leg
$-0.49\%$, long leg $+0.46\%$ a month relative to F\&O stocks), which is a small-stock effect rather than the one-sided
pattern of a short-sale constraint. Momentum alone fits the prediction: losers among constrained stocks do worse ($-0.51\%$
a month, s.e.\ 0.22). The groups differ greatly in size (median market capitalisation Rs 43,513 crore for F\&O stocks
against Rs 770 crore for cash-only stocks without SLB trades), which limits what this comparison can show
(Figure~\ref{fig:h5}).

\begin{figure}[t]
\caption[Anomaly returns where shorting is hardest]{Value-weighted, the overpriced side does no worse where shorting is hardest}\label{fig:h5}
\exsub{Most constrained group (cash-only stocks without SLB trades) minus F\&O stocks, short and long legs of a composite mispricing sort}
\centering\includegraphics[width=\linewidth]{fig_h5_legs}
\exsource{Newey-West standard errors; holding months January 2016 to September 2026. Source: author's calculations.}
\end{figure}

\section{Drift and price delay across all stocks}
Two further comparisons use the whole cash market. First, relative to F\&O stocks of the same size decile, cash-only stocks
drift $-0.67\%$ more after bad news (s.e.\ 0.42) and $+2.40\%$ more after good news (s.e.\ 0.43): they under-react in both
directions. Second, the constraint index of Chapter~\ref{ch:data}, which combines the lack of a short channel with SLB
scarcity, spreads, illiquidity and idiosyncratic volatility, predicts slower price adjustment the next year, also within a
stock over time and after controlling for size. It also predicts more drift after bad news than after good news, but that
asymmetry disappears once size is controlled for, in line with the symmetric under-reaction of cash-only stocks.

\section{Summary}
Short-sale access matters for how fast prices adjust: entry into the F\&O segment speeds the response to market news, and
more constrained stocks adjust more slowly. The evidence on overpricing and on drift after bad news is mixed: entry does
not weaken bad-news drift, value-weighted anomaly returns do not show the predicted asymmetry, and the link between the
broader constraint index and drift after bad news disappears once size is controlled for.
